
\needspace{4\baselineskip}
\color{uniblue}{\subsection{Предупредительная сигнализация устройства (ПС)\label{ps:punkt}}}
\color{black}

\begin{enumerate}[label=\arabic{section}.\arabic{subsection}.\arabic*, labelsep=4pt, leftmargin=0pt, itemindent=57pt]

\item	
Предупредительная сигнализация предназначена для информирования оперативного персонала при отклонениях от нормального режима работы оборудования энергообъекта.

\item
В устройстве предупредительная сигнализация представлена функцией «ПС».

\item
Функция принимает сигналы срабатывания защитных и сервисных функций устройства и формирует следующие сигналы:
\begin{itemize}
	\item импульсный сигнал предупредительной сигнализации «ПС: Пуск (имп)»;
	\item длительный сигнал предупредительной сигнализации «ПС: Пуск».
\end{itemize}

\item
Выходные сигналы функции могут быть использованы для дальнейшего конфигурирования (например, на выходные реле) и применяться в цепях участковых и индивидуальных табло, сигнальных ламп, сигнальных реле, звуковой сигнализации и пр.

\item
Введенный ключ «Вывод терминала» блокирует формирование выходных сигналов.

\item
Сигналы из логики СС: «ТЗ сигн», «ГЗ сигн», «Низк.изол.ТЗ», «Низк.изол.ГЗ», «ТЗ заблокирована», «ГЗ заблокирована», «Вых.цепи разобраны», «БИ выведены», «ОТ сигн», «Неисправность КА», «Прев. вр.пер. КА», «Общ.внеш.сигнал» контролируются соответствующими уставками SGF1\ldots SGF12.

\item
Инвертированное состояние дискретного сигнала при введенной уставке SGF13 «Полож.пер.ав.дебл.» формирует работу логики ПС.

\item
Логика работы функции выполнена в соответствии с рисунком приложения \ref{app:ps}. В таблице \ref{ps:tbl1} приведены параметры, необходимые для настройки функции. 

\small
\begin{longtable}{|>{\centering\arraybackslash}m{5.3cm}|>{\centering\arraybackslash}m{3.3cm}|>{\centering\arraybackslash}m{4.2cm}|>{\centering\arraybackslash}m{1.8cm}|>{\centering\arraybackslash}m{1cm}|}
\caption{Параметры для настройки предупредительной сигнализации\hfill\vspace{-0.5\baselineskip}}\label{ps:tbl1}\\ 
\hline
\rowcolor{gray!30}
Параметр (Параметр на ИЧМ) & Условное обозначение на схеме & Значение/ Диапазон & Единица измерения & Шаг \\ 
\hline
\endfirsthead
%\caption{\emph{Продолжение\hfill\vspace{-0.5\baselineskip}}} \\ % Отступ обозначения таблицы от таблицы и выравнивание надписи слева
\caption*{\hspace{3pt}\emph{Продолжение таблицы \ref{ps:tbl1}\hfill\vspace{-0.5\baselineskip}}} \\ % сделано по ГОСТ 2.105 п.6.8.7
\hline
\rowcolor{gray!30}
Параметр (Параметр на ИЧМ) & Условное обозначение на схеме & Значение/ Диапазон & Единица измерения & Шаг \\ 
%\hline
\endhead
%\hline
\endfoot
\hline
\endlastfoot
%==+t1*CALH1|LVALH
\centering Контроль сигнализации ТЗ (Контр\_ТЗ\_сигн) & \centering SGF1 & \centering 0 = Не предусмотрено\\1 = Предусмотрено & \centering -- & \centering \arraybackslash -- \\
\hline
\centering Контроль сигнализации ГЗ (Контр\_ГЗ\_сигн) & \centering SGF2 & \centering 0 = Не предусмотрено\\1 = Предусмотрено & \centering -- & \centering \arraybackslash -- \\
\hline
\centering Контроль низкой изоляции ТЗ (Контр\_низк\_изол\_ТЗ) & \centering SGF3 & \centering 0 = Не предусмотрено\\1 = Предусмотрено & \centering -- & \centering \arraybackslash -- \\
\hline
\centering Контроль низкой изоляции ГЗ (Контр\_низк\_изол\_ГЗ) & \centering SGF4 & \centering 0 = Не предусмотрено\\1 = Предусмотрено & \centering -- & \centering \arraybackslash -- \\
\hline
\centering Контроль блокировки ТЗ (Контр\_ТЗ\_блок) & \centering SGF5 & \centering 0 = Не предусмотрено\\1 = Предусмотрено & \centering -- & \centering \arraybackslash -- \\
\hline
\centering Контроль блокировки ГЗ (Контр\_ГЗ\_блок) & \centering SGF6 & \centering 0 = Не предусмотрено\\1 = Предусмотрено & \centering -- & \centering \arraybackslash -- \\
\hline
\centering Контроль выходных цепей (Контр\_цепей) & \centering SGF7 & \centering 0 = Не предусмотрено\\1 = Предусмотрено & \centering -- & \centering \arraybackslash -- \\
\hline
\centering Контроль выведенного положения БИ (Контр\_вывода\_БИ) & \centering SGF8 & \centering 0 = Не предусмотрено\\1 = Предусмотрено & \centering -- & \centering \arraybackslash -- \\
\hline
\centering Контроль сигнализации опер.тока (Контр\_ОТ\_сигн) & \centering SGF9 & \centering 0 = Не предусмотрено\\1 = Предусмотрено & \centering -- & \centering \arraybackslash -- \\
\hline
\centering Контроль неисправности КА (Контр\_Неиспр\_КА) & \centering SGF10 & \centering 0 = Не предусмотрено\\1 = Предусмотрено & \centering -- & \centering \arraybackslash -- \\
\hline
\centering Контроль превышения времени переключения КА (Контр\_прев\_вр\_пер) & \centering SGF11 & \centering 0 = Не предусмотрено\\1 = Предусмотрено & \centering -- & \centering \arraybackslash -- \\
\hline
\centering Контроль общего внешнего сигнала (Контр\_общ\_вн\_сигн) & \centering SGF12 & \centering 0 = Не предусмотрено\\1 = Предусмотрено & \centering -- & \centering \arraybackslash -- \\
\hline
\centering Контроль положения переключателя аварийной деблокировки  (Контр\_пол\_ав\_дебл) & \centering SGF13 & \centering 0 = Не предусмотрено\\1 = Предусмотрено & \centering -- & \centering \arraybackslash -- \\
%===t1
\end{longtable}
\normalsize

\end{enumerate}